\chapter{Create the Test Runner Jail}

The application stack is now managed independently of any terminal. We can add
one final jail whose only job is to observe the system through nginx. The test
runner will know the web URL and expected page content, but nothing about the
store controller, catalog API, or database schema.

This runner enters through nginx's private address. Chapter 17 separately
proved the host PF boundary from another LAN machine; keeping the automated
runner on the managed network makes the test repeatable without depending on a
particular host LAN address or forwarded host port.

\section{Exercise 21A: Create the blank test jail}

Open BSDRunner Jails and choose \textbf{Create a Service Jail}. Use:

\begin{description}
  \item[starter] Blank FreeBSD;
  \item[jail name] \texttt{test};
  \item[FreeBSD release] \texttt{15.1-RELEASE};
  \item[private address] \texttt{10.17.89.27/24};
  \item[outbound policy] LAN + Internet;
  \item[starts at boot] disabled; and
  \item[inbound exposure] none.
\end{description}

The broader outbound policy is deliberate. Internet access lets the jail
install the Ruby test toolchain; private-network access lets it reach nginx at
\texttt{10.17.89.26}. It still receives no inbound host port.

Review the plan, type \texttt{test} into the confirmation field, and create the
jail. Wait until it is running.

\section{Exercise 21B: Probe the black-box boundary}

Open the test jail console and request the page using only FreeBSD's base
tools:

\begin{lstlisting}[language=sh]
fetch -qo - http://10.17.89.26/ | \
  grep 'FreeBSD Handbook'
\end{lstlisting}

Expect the rendered product line. This request originates outside all three
service jails and crosses nginx, store, catalog, and PostgreSQL.

\section{Exercise 21C: Declare the test dependency}

Open the \texttt{test} jail's \textbf{Manifest} editor. Add this line with its
top-level jail settings:

\begin{lstlisting}
depends_on = ["web"]
\end{lstlisting}

Run \textbf{Dry run}, review the dependency change, then choose
\textbf{Apply}. The complete declared graph is now:

\begin{lstlisting}
test -> web -> store -> catalog -> database
\end{lstlisting}

Bringing up \texttt{test} can therefore bring up every missing dependency in
order. The test itself remains an ordinary client request rather than gaining
special access to those services.

\section{Checkpoint}

Stop after the request from the test jail finds \texttt{FreeBSD Handbook} and
the manifest records its dependency on web. Leave the console open. The next
exercise will install Ruby and express the same observation as executable
Minitest assertions.
